
\begin{refsection}

\chapter*{Introducción}
\addcontentsline{toc}{chapter}{Introducción}
\label{cap-21}

La pandemia fue un acontecimiento disruptivo, que lo transformó todo durante un tiempo considerable y dejó huellas que aún no terminamos de vislumbrar. A seis años del comienzo de la crisis pandémica global y de las medidas de aislamiento en la Argentina, en este libro presentamos una reconstrucción de las experiencias de la vida cotidiana durante ese tiempo diferente. Nuestro objetivo fue visibilizar y dar cuenta de las distintas formas en que se atravesó la pandemia. Aunque fue un hecho social total, produjo una heterogeneidad de experiencias, en parte siguiendo los patrones conocidos de estructuración social como el género, la clase y la edad, pero también generando nuevas desigualdades.

Como investigadores e investigadoras, la pandemia nos atravesó. No solo en términos personales, en nuestras rutinas y organización familiar, sino que también cambió nuestras agendas de trabajo, la forma de dar clases y de interactuar con colegas. En particular para quienes analizamos las vidas cotidianas, las prácticas culturales y las tecnologías digitales, fue un momento desafiante de rápidas transformaciones y a la vez, en términos del trabajo intelectual, fue una etapa de sinergias que nos permitió pensar juntos un problema en sus diferentes aristas.

El libro presenta los resultados de un proyecto de investigación\footnote{PIP CONICET «Pandemia y vida cotidiana en el AMBA: un abordaje
    de la heterogeneidad de las experiencias», con sede en la
    Escuela Interdisciplinaria de Altos Estudios Sociales, Universidad
    Nacional de San Martín. Dirigido por Marina Moguillansky e integrado
    por Carolina Duek, Sebastián Benítez Larghi, Magdalena Lemus, Mariana
    Lopresti, Rosario Guzzo, Julián Kopp, Paula Simonetti y Pablo Salas.} que se propuso explorar, desde la sociología de la cultura, las transformaciones de la vida cotidiana a raíz de la pandemia en Buenos Aires y alrededores, en un contexto signado por desigualdades. Para ello, conversamos con diferentes personas sobre sus experiencias de la pandemia, los cambios en sus rutinas laborales, en la educación, en sus prácticas de sociabilidad, los consumos culturales y en el uso de las pantallas. En este análisis, partimos de la hipótesis de que hubo muchas y muy diversas experiencias de la pandemia, pero que a su vez, esa heterogeneidad se puede explicar y comprender mejor a partir de considerarla a la luz de las posiciones de clase, género, generación, la situación familiar y el lugar de residencia. Nos propusimos centrar la indagación en la dimensión subjetiva y emocional de la pandemia, las formas personales de valorar lo vivido y los horizontes a futuro. Por eso privilegiamos un enfoque biográfico y narrativo que diera lugar al relato de la experiencia singular y a los recorridos de cada persona. Las historias, las experiencias y los recuerdos se cruzaron en los relatos de distintas formas con un eje común: lo vivido fue excepcional y tuvo un gran impacto en todos los ámbitos de la vida social, familiar e individual.

«¿Fue en 2020 o en 2021? No recuerdo bien», se repitió en muchas entrevistas como síntesis de la ruptura de la cadencia esperable del devenir del tiempo, de las rutinas y de la cotidianidad. Cuando indagamos sobre el regreso de ciertas actividades, la vacunación y el fin de las restricciones, muchos relatos confundían los años los tiempos y las medidas de aislamiento dando cuenta de que la línea temporal se rompió en muchos pedazos y su reconstrucción era aún frágil e imprecisa. Este es uno de los tantos ejemplos que ilustra la disrupción que supuso la pandemia como hecho social total. En cada uno de los capítulos analizaremos historias, trayectorias y experiencias biográficas que, en conjunto, nos permiten ensayar una reconstrucción parcial de las memorias, las huellas y los cambios que dejó la pandemia.


\section{¿Qué sabemos sobre la pandemia y sus efectos en la vida cotidiana?}
\label{cap-21.que-sabemos-sobre-la-pandemia-y-sus-efectos-en-la-vida-cotidiana}
Una de las investigaciones de mayor alcance y profundidad sobre la pandemia y sus efectos sociales fue coordinada por Pablo Dalle \protect\parencite*{4181-DALLE2022}. Analizaron datos de la Encuesta Permanente de Hogares (EPH) y del Sistema Integrado Previsional Argentino (SIPA) e hicieron una encuesta a nivel nacional, entrevistas y estudios de caso cualitativos, buscando explorar los efectos de la doble crisis económica y pandémica sobre la estructura social. Afirman que Argentina atravesó una doble crisis: por un lado, una previa que se inicia en 2016 con el cambio de gobierno y las políticas de endeudamiento, ajuste del gasto y aumento de la inflación; a esta crisis se suma el impacto de la pandemia en 2020 con una fuerte caída de la actividad económica y del empleo. La hipótesis de la investigación es que la pandemia y el aislamiento tuvieron efectos diferenciales en los individuos y las familias, de acuerdo a sus posiciones de clase. Así, los diversos capítulos muestran el papel de la clase social en las respuestas a la pandemia por parte de las familias y los hogares; se trata del factor principal para explicar las probabilidades de perder el trabajo; explica la diversidad de estrategias habitacionales desarrolladas y también la evaluación de las políticas del Estado para mitigar los efectos sociales del aislamiento. El impacto de la pandemia fue mayor sobre trabajadores informales y pequeña burguesía, que se corresponde con las clases medias bajas; produjo una pérdida de bienestar en los hogares de la base de la pirámide social. En este sentido, las desigualdades con las que trabajaremos en los diferentes capítulos se articulan con las estrategias de los sujetos sociales en sus diferentes posiciones relativas en la estructura social.

En Argentina, la gestión de la pandemia tuvo una gravitación central y se convirtió en organizadora de las rutinas de la vida cotidiana. Con mayor o menor impacto según las regiones (aquí nos ocupamos de la región metropolitana, es decir, la ciudad de Buenos Aires y sus alrededores) y según el tipo de actividad de las personas. El Estado definió quiénes podían circular y quiénes no, implementó políticas de mitigación de la crisis y reguló la apertura de espacios y actividades, modalidades de vinculación en las familias con padres y madres separados entre otras cuestiones. Mariana Heredia y su equipo indagaron sobre el papel del Estado argentino en la pandemia en el libro \emph{Lo que pudo (y lo que no pudo) el Estado}. Allí analizan las principales políticas: el Ingreso Familiar de Emergencia (IFE) y el programa de Asistencia de Emergencia al Trabajo y la Producción (ATP). El diseño de las políticas, sus formas de implementación, los mediadores técnicos, la experiencia de empresas y beneficiarios, los problemas de gestión y a los que se enfrentaron fueron inéditos por la premura y el alcance de estas medidas.

La investigación de Heredia y equipo puso de relieve la gran heterogeneidad de experiencias de la pandemia que es el punto de partida de nuestra propia investigación. Dada la diversa composición social y económica de las regiones y provincias de la Argentina, «el riesgo sería creer que hubo una sola experiencia de la pandemia y de las políticas de asistencia»
         \protect\parencite[p. 15]{4182-HEREDIA2022}. Por ejemplo, el IFE tuvo una gran importancia en provincias del Norte del país, llegando a alrededor del 40\% de la población, mientras que en CABA alcanzó a menos del 20\%; en cuanto al ATP, por el contrario, observan que alcanzó a cerca del 70\% de las empresas privadas radicadas en CABA y alrededor del 50\% de las de la provincia de Buenos Aires, con porcentajes mucho menores en otras regiones. Además, los diferentes sectores de actividad fueron afectados en mayor o menor medida y «experimentaron la crisis de manera heterogénea»
         \protect\parencite[p. 172]{4178-ALVAREZ2022}, de modo que la inserción laboral de los sujetos y el sector de su actividad fue un condicionante fuerte de sus experiencias de la pandemia. Entre los sectores más afectados estuvieron la hotelería y restaurantes, los servicios personales y la construcción; también tuvieron grandes caídas la manufactura, el comercio y las empresas de transporte; mientras que los que menos cayeron en sus niveles de actividad fueron la agricultura y ganadería, los servicios públicos (gas, luz, agua) y los servicios financieros. El impacto demográfico de la destrucción de puestos de trabajo fue mayor en las mujeres y en los jóvenes. Algunas familias de clase media-baja fueron por primera vez objeto de políticas de asistencia social durante la pandemia y que ello supuso un posicionamiento subjetivo novedoso y problemático (y también da cuenta de cómo las desigualdades se ampliaron llegando a nuevos espacios del entramado social).

En una clave cercana a nuestro trabajo, el \emph{Grupo de Estudios Culturales y Urbanos} que coordina Juliana Marcus, estudió la experiencia urbana de la pandemia a través de una combinación de dos encuestas online, entrevistas y observaciones en espacios públicos como parques y plazas. En el libro \emph{Ciudad confinada. Experiencias urbanas durante la pandemia covid-19}
         \protect\parencite*{4183-MARCUS2023} Marcus \emph{et al.}, señalan que la pandemia y el confinamiento trastocaron las rutinas de la vida cotidiana, en particular el vínculo con la ciudad; diversas prácticas se relocalizaron al interior de las viviendas (el trabajo, la educación, el ocio y la sociabilidad) y también, cuando se flexibilizaron las restricciones, aparecieron usos novedosos de los espacios públicos (deportes en grupo, festejos de cumpleaños y reuniones sociales se trasladaron a los espacios verdes como parques o plazas). Entre sus resultados, destacan el pasaje desde el miedo y la evitación de los espacios públicos urbanos hacia una revalorización y reencantamiento de los mismos. De la urbanofobia a la urbanofilia es la fórmula que condensa este pasaje. Si en los primeros meses de pandemia se registró un miedo exacerbado o fobia social a salir del hogar, esto fue promovido por discursos y políticas estatales tanto a nivel nacional como provincial y local. Con el paso del tiempo, las medidas de aislamiento se flexibilizaron y las autoridades empezaron a promover las actividades en espacios abiertos. Otro de los hallazgos clave de Marcus \emph{et al.} es que la configuración familiar y/o del hogar fue un factor central para la experiencia del aislamiento. Vivir solo, estar en pareja, tener a cargo hijos pequeños o adultos mayores modularon los recursos y las posibilidades para sostener o no el trabajo, para reestructurar las rutinas de ocio y sociabilidad, y condicionaron los niveles de control y autonomía sobre el propio tiempo.

El equipo dirigido por Javier Balsa analizó cómo la pandemia generó cierta polifonía discursiva en la que se entrecruzaban y a veces tensionaban las narrativas oficiales, las de la oposición política y las de medios de comunicación. Este análisis subraya la coexistencia de discursos que, por un lado, buscaban legitimar medidas de control y, por otro, expresaban críticas y demandas. A través de grupos focales, entrevistas y encuestas telefónicas, se indagó cómo eran percibidas y narradas las medidas de confinamiento. La perspectiva de Balsa y equipo \protect\parencite*{4177-BALSA2023} destaca el papel de ciertas metáforas –como la bélica– en los discursos oficiales y mediáticos, que estructuraron e instalaron sentidos sobre la percepción de la emergencia. Estas construcciones simbólicas jugaron un rol importante en la forma en que se interpretaron las medidas restrictivas y en la movilización de la opinión pública.

Por otro lado, el aislamiento en el hogar, la regulación de las fronteras y de la circulación de sujetos durante los primeros meses de la pandemia (y un poco más también), generaron modos novedosos de imaginar el espacio de la vivienda, sus entornos y la ubicación geográfica. Este fue el eje del trabajo de investigación que dirigió Ramiro Segura \protect\parencite*{4188-SEGURA2024}, con una amplia red de equipos en distintos espacios del país. Las restricciones a la circulación y las medidas de prevención de contagio multiplicaron las barreras y regulaciones sobre la movilidad cotidiana, e instalaron en el discurso público la idea del AMBA.

Todas estas investigaciones fueron puntos de partida y sus responsables son interlocutores necesarios del análisis que presentamos en este libro. En conjunto plantean que la pandemia fue un momento disruptivo, que hubo formas diferentes de atravesarla y que tuvo consecuencias también diferenciales según la configuración familiar, la clase social, el género y la edad. A ello agregamos como clave interpretativa la atención a los recursos subjetivos y afectivos para transitar esta experiencia y para su posterior elaboración simbólica.


\section{El método biográfico, la vida cotidiana y las experiencias de la pandemia}
\label{cap-21.el-metodo-biografico-la-vida-cotidiana-y-las-experiencias-de-la-pandemia}
Nuestro análisis de las experiencias de la pandemia parte de considerar la dimensión subjetiva de las transformaciones de la vida cotidiana, las formas en que se experimentaron los cambios y las estrategias familiares o personales que se llevaron adelante durante ese período. La metodología de la investigación partió del enfoque interpretativo y biográfico, como herramientas para escudriñar las tramas de la vida cotidiana.

Según Agnes Heller, la vida cotidiana es el conjunto de las actividades de reproducción de los hombres; no es abstracta sino concreta –particular, específica– y no solo tiene una historia sino que se la puede considerar a su vez como «espejo de la historia»
         \protect\parencite[p. 20]{4166-HELLER1977}. Recuperando ideas de Marx, Heller escribe que toda revolución social cambia radicalmente la vida cotidiana; pero que con la distancia del tiempo, el historiador descubre que en ella puede encontrar que las transformaciones se anticipan en las estructuras de lo cotidiano. Esta dialéctica entre el modo de vida y las transformaciones quizás nos sirva para pensar cómo la pandemia radicalizó algunas formas de interacción social que ya venían instalándose desde tiempo atrás. Entre ellas, la creciente reclusión hogareña, que ya había tematizado Roman Gubern en \emph{El simio informatizado}
         \protect\parencite*{4184-GUBERN1987} al referirse a la «cueva aterciopelada» que resulta de la ampliación de bienes y servicios domésticos y la percepción del afuera como un espacio hostil y peligroso. La multiplicación de pantallas y el desarrollo del entorno digital también es un proceso previo y que las ciencias sociales tematizaron al menos desde \emph{La era de la información}
         \protect\parencite{4179-CASTELLS1998}, como un cambio cultural que lleva varias décadas pero que a su vez se intensificó en los últimos años, dando lugar a una «mediatización profunda», como muestran Boczkowski y Mitchelstein en \emph{El entorno digital}
         \protect\parencite*{4187-MITCHELSTEIN2022}. En este sentido, las condiciones materiales para trasladar actividades laborales, educativas y de ocio al hogar ya estaban presentes antes de la pandemia. Pero, nuevamente, no todas las familias ni personas que atravesaron la pandemia tuvieron las mismas posibilidades en lo que respecta a la conectividad y a las modalidades de interacción mediante plataformas virtuales.

El camino que elegimos para analizar relatos de la vida cotidiana durante la pandemia y reconstrucciones de la memoria de ese período se ordenaron en torno del método biográfico. Siguiendo los desarrollos claves de Arfuch \protect\parencite*[p. 22]{4175-ARFUCH2002}, la idea de \emph{espacio biográfico} es muy relevante para dar cuenta de «la multiplicidad, lugar de confluencia y de circulación, de parecidos de familia, de vecindades y de diferencias». El interés en la voz y en la experiencia de los sujetos con un énfasis en la dimensión testimonial, según Arfuch, constituye un símbolo del siglo XXI que se ha obsesionado por los hitos de la memoria. Según su hipótesis y a los fines de este libro y de nuestros análisis, plantea que lo biográfico es un espacio intermedio entre lo público y lo privado. No se trata solamente de lo que recordamos sino de cómo lo que nos pasó de manera directa o indirecta se puede relacionar con la dimensión pública y experiencial de un momento en el tiempo y en la historia.

Nos inspiramos también en la perspectiva de la sociología pública feminista que propone Kate Averett \protect\parencite*[p. 1]{4176-AVERETT2021}, que resalta la importancia ética de dar lugar a la elaboración narrativa del sufrimiento durante la pandemia. Según su planteo, es relevante «dar testimonio de las experiencias de los afectados por la pandemia, hacer un registro de esas experiencias y ayudar a los académicos y al público a pensar sociológicamente sobre la pandemia, idealmente de tal manera que ayude en la creación de políticas respuestas que aborden y reduzcan este sufrimiento». La propia investigación sociológica que se desarrolla a través de entrevistas en profundidad, en las cuales los sujetos nos cuentan y se cuentan a sí mismos sus vivencias, tiene un carácter performativo al ofrecer un espacio de escucha atenta, de catarsis y elaboración simbólica, diferente de la conversación con pares o de la escucha terapéutica.

Como señala Ernesto Meccia en \emph{Biografía y sociedad}, los métodos biográficos pueden orientarse a reconstruir hechos o experiencias; en nuestro caso, optamos por la segunda alternativa. El diseño metodológico del proyecto se basó en la realización de entrevistas cualitativas de carácter biográfico, centradas en el recorrido y la historia de vida de cada persona hasta llegar a la pandemia, con la intención de reconstruir sus experiencias, esto es, «las formas que tiene la gente de significar esos hechos por intermedio de su propia memoria biográfica»
         \protect\parencite[p. 25]{4186-MECCIA2022} Cómo estaban cuando comenzó el aislamiento, de qué maneras cambiaron sus rutinas cotidianas y laborales, cómo se modificó (o no) la configuración de sus hogares y familias, cómo recuerdan este período, son algunos de los interrogantes que guiaron las entrevistas.

Para la selección de los entrevistados, utilizamos la lógica del muestreo teórico en etapas sucesivas \protect\parencite{4167-RAPLEY2014}, incorporando en cada ronda los hallazgos de los casos ya analizados como nuevos criterios de selección. Esta modalidad de trabajo supone un análisis simultáneo a la realización de las entrevistas, para poder de esta manera obtener orientaciones que surgen de los propios datos y definir las estrategias de muestreo, hasta lograr la saturación teórica \protect\parencite{4180-CHARMAZ2006}. Así, como el género aparecía como relevante desde el comienzo, en las rondas de selección de entrevistados buscamos mantener la representación de hombres y mujeres; la edad y la situación vital también eran configuradores de experiencias diferenciales, por lo cual buscamos entrevistar adolescentes, jóvenes, adultos y ancianos. La posición de clase, el tipo de trabajo realizado y el sector de actividad también fueron variables importantes. La identidad política\footnote{Una primera forma de aproximarnos a la identidad política para
      seleccionar a los entrevistados fue considerar a quién habían votado
      en las elecciones presidenciales de 2019, en primera y segunda vuelta,
      lo que nos permitía dividir al universo en dos grandes grupos. Sin
      embargo, esto se reveló incongruente con las posiciones políticas de
      los entrevistados que fuimos incorporando, mostrando signos del
      creciente desencanto de una parte significativa de la población.} apareció desde el comienzo como factor clave que incidía en la percepción subjetiva del riesgo de enfermarse, en el cumplimiento de las medidas de restricción y en la evaluación del rol del gobierno. Si bien al comienzo hubo consenso político en la gestión de la pandemia, con el tiempo surgieron conflictos y cuestionamientos (Spólita, Balsa y Brusco, 2022). El lugar de residencia (en CABA o AMBA) también implicaba diferencias en cuanto a cómo se experimentaban las restricciones. Por último, ser considerado o no como trabajador «esencial» fue un diferenciador de la experiencia del aislamiento y también fue un criterio de selección. Así, el objetivo fue lograr la máxima variabilidad en el perfil de los informantes.

Las entrevistas fueron realizadas entre enero y julio de 2022. Dado el contexto de cercanía con la pandemia, alrededor de la mitad de esas entrevistas fueron realizadas en forma virtual, con videollamadas. Entrevistamos a setenta y ocho personas, de las cuales 40 fueron varones y 38 fueron mujeres. Las edades de los entrevistados van desde los 14 hasta los 86 años, con mayor representación del grupo de adultos (entre 30 y 60 años).

Para construir la clasificación socioeconómica de las personas entrevistadas, partimos de las preguntas del esquema SAIMO (Sociedad Argentina de Investigadores de Marketing), que agrupan a la población en cinco niveles según el nivel educativo, el tipo de inserción laboral y el acceso a prestaciones sociales \protect\parencite{4185-CLEMENCEAU2016}. Esta segmentación inicial nos permitió organizar las rondas de muestreo para cubrir todas las posiciones de la estructura social. Pero además de esos datos tipificados, para cada entrevistado contamos con información detallada sobre su lugar de residencia, la trayectoria formativa, el tipo de empleo, la composición familiar, el acceso a bienes y prácticas culturales, lo cual nos permitió matizar y validar la clasificación de cada sujeto. A partir de estos datos, en los capítulos analíticos empleamos categorías conceptuales propias de las ciencias sociales –como clase social, condiciones de vida y trayectorias laborales– que facilitan una comprensión más situada y compleja de las desigualdades y vivencias durante la pandemia. Así, para describir la composición social del conjunto de la muestra, podemos decir que entrevistamos a doce dueños de pequeñas y medianas empresas y profesionales en ocupaciones de alta remuneración; a cuarenta y tres personas de posiciones medias y de clase media-baja, incluyendo a empleados públicos y privados, profesionales independientes, emprendedores, docentes, entre otros; y a veintidós personas de bajos o muy bajos ingresos, de sectores populares, comprendiendo a desocupados, trabajadores informales en sectores marginales, cartoneros, jubilados, cuidadores de ancianos y otros perfiles similares.

El análisis de las entrevistas se hizo de manera colaborativa, con la codificación de las transcripciones y la discusión orientada por interrogantes y conceptos sensibilizadores: rutinas de la vida cotidiana, organización del trabajo, control, autonomía o dependencia en la adopción de estos cambios, percepción del riesgo. La sociología de las emociones fue también una herramienta importante para la escucha y el análisis de las entrevistas: en los relatos aparecía con insistencia el miedo, la angustia, la tristeza; en otros casos el enojo, la frustración y la impotencia. Para ello, resultó fundamental incorporar al análisis las diferentes formas de la comunicación no verbal como la entonación, las pausas, los gestos con el rostro y las manos, la mirada, la risa y el llanto.


\section{El recorrido de este libro}
\label{cap-21.el-recorrido-de-este-libro}
Este libro reúne una serie de capítulos que abordan, desde distintas perspectivas, las experiencias cotidianas durante la pandemia de COVID-19 en el Área Metropolitana de Buenos Aires (AMBA). El hilo conductor que organiza el volumen es la dimensión generacional: los textos se ordenan siguiendo un recorrido que va desde las vivencias de adolescentes hasta las de personas mayores. Este criterio permite reconstruir cómo se vivió el aislamiento social en distintas etapas del ciclo vital, y cómo la edad se entrelazó con otros factores —como el género, la clase social, las condiciones habitacionales o el tipo de inserción laboral— para configurar experiencias desiguales del encierro. Aunque el orden sugiere una progresión cronológica, los capítulos pueden leerse también de forma autónoma, ya que cada uno presenta casos particulares y ofrece claves analíticas propias.

El libro se abre con el capítulo de Mariana Lopresti, «Pantallas apagadas y una difícil convivencia», centrado en las experiencias de adolescentes durante la suspensión de clases presenciales. A partir de cuatro historias de jóvenes de sectores medios del AMBA, se analizan las formas de habitar el encierro, las tensiones familiares derivadas de la convivencia intensificada y la importancia de la escuela como espacio de sociabilidad en esta etapa de la vida. Las historias de los y las adolescentes entrevistados dan cuenta de la ausencia de contacto cotidiano con los amigos, que muchas veces derivó en la ruptura de lazos, pero también los cambios rápidos en el cuerpo y la sensación de haber perdido para siempre ciertas experiencias vitales. El enfoque destaca cómo el género y la edad condicionaron las prácticas y sentidos asociados al espacio doméstico, y pone de relieve la capacidad de los y las adolescentes para reflexionar críticamente sobre su experiencia.

El siguiente capítulo es de Magdalena Lemus, titulado «Trabajar, estudiar y cuidar en pandemia». Lemus reconstruye las trayectorias de tres jóvenes de distintas clases sociales que atravesaron el aislamiento tratando de sostener sus estudios, sus lazos sociales y, en algunos casos, sus trabajos. El análisis muestra cómo las condiciones materiales —como el acceso a tecnologías o la necesidad de generar ingresos— marcaron diferencias significativas entre los casos, pero también señala puntos de encuentro: la centralidad de la sociabilidad entre pares y el protagonismo del cuidado de otras personas como dimensión clave de la experiencia juvenil durante la pandemia.

En «Aislados y desorganizados», Sebastián Benítez Larghi se enfoca en familias de sectores populares que perdieron su empleo o vieron interrumpidos sus ingresos. A través de dos historias de vida, el capítulo muestra cómo se desplegaron y reforzaron los mecanismos de reproducción de las desigualdades sociales, especialmente a partir del debilitamiento de los vínculos sociales durante la pandemia. La vivencia de la pandemia en estos hogares estuvo marcada por la desestructuración y resignificación de unas condiciones de existencia que ya eran frágiles y precarias. Los escasos puntos de apoyo que sostenían su vida cotidiana fueron rápidamente arrasados por una crisis que desbordó las relaciones entre las personas, las instituciones y los conocidos modos de integración social: pérdida de empleo, interrupción de ingresos, incomprensión de las políticas de asistencia estatal, dificultades para garantizar la subsistencia, desconexión del entramado educativo y falta de respuesta de los sistemas de salud.

Carolina Duek, en «Todos juntos, todo el día, todos los días», analiza el modo en que tres mujeres con trabajo previo a la pandemia reorganizaron sus rutinas y espacios domésticos para sostener sus ingresos durante el aislamiento. A través de sus relatos, el capítulo muestra cómo la pandemia alteró la experiencia de esas mujeres que trabajaron desde el hogar y que transformaron sus formas de habitar la casa, generando nuevas estrategias para poder trabajar y a su vez nuevas formas de lidiar con las dimensiones afectivas del aislamiento. El texto propone un análisis centrado en los núcleos comunes entre situaciones diversas y la necesidad de repensar el hogar como espacio de producción, cuidado y tensión.

El capítulo de Marina Moguillansky, «¿Todo sigue igual?», aborda la experiencia de los trabajadores esenciales, a través de tres historias de hombres empleados respectivamente en el sector de salud, la recolección de residuos y la producción de alimentos. Si bien estos trabajadores siguieron saliendo a la calle, sus rutinas se transformaron profundamente. A partir de tres relatos, se pone en evidencia la paradoja de una continuidad que es, al mismo tiempo, una ruptura: las calles vacías, los nuevos protocolos, el contacto directo con la enfermedad y la muerte. El capítulo ofrece una mirada sobre las formas de habitar una ciudad desierta y de enfrentar riesgos laborales en condiciones extremas.

Por último, Rosario Guzzo, en su capítulo sobre «Cuidar y descuidar», analiza las experiencias de personas mayores que fueron incluidas en la categoría de grupo de riesgo que se consideraban más vulnerables al contagio. A través de tres trayectorias biográficas, indaga en las tensiones entre el aumento del cuidado recibido y la pérdida de autonomía, así como en el uso y apropiación de tecnologías digitales para sostener ciertas rutinas. El texto resalta que, si bien se reconoció la necesidad de mayor cuidado, este fue absorbido principalmente por las familias, lo cual subraya la urgencia de políticas de cuidado que preserven la autonomía de las personas mayores.

En su conjunto, el libro muestra cómo la pandemia alteró las coordenadas de la vida cotidiana, pero también cómo esas alteraciones se vivieron de manera desigual. Algunos relatos comparten preocupaciones en torno a la reorganización del espacio doméstico, el sostenimiento de los vínculos o las estrategias de adaptación, mientras que otros muestran rupturas drásticas, marcadas por la precariedad, la soledad o el riesgo. Las experiencias varían, pero también dialogan entre sí, trazando un mapa complejo de lo vivido durante el aislamiento. Este libro no busca ofrecer una síntesis totalizadora, sino reconstruir desde la particularidad y la comparación, desde las biografías y los contextos, algunas formas de habitar lo cotidiano cuando lo cotidiano se volvió extraordinario. A través de múltiples voces, invita a pensar la pandemia como una experiencia generacional y social, atravesada por desigualdades estructurales, pero también por estrategias de sentido, cuidado y resistencia.


\section*{\refname}
\printbibliography[heading=none]
\end{refsection}

